\section*{Chapter \ref{ch:ml:validation} --- Validation}

\begin{solution}{exo:ml:validation:1}
$0.30\% + 1.16\times0.10\% = 0.42\%$: the winner's score overstates every candidate's true 0.30\% by 0.12 points.
\end{solution}

\begin{solution}{exo:ml:validation:2}
The test labels span months 100 to 141. Purging drops the training labels starting at 98 and 99 (they end inside the
span) and at 140 and 141 (they start inside it); the embargo also drops 142.
\end{solution}

\begin{solution}{exo:ml:validation:3}
Mean 0.18\%, standard deviation 0.03\%, standard error $0.03/\sqrt3 = 0.017\%$: ten standard errors above zero on
targets that cannot be predicted. The pipeline leaks.
\end{solution}

\begin{solution}{exo:ml:validation:4}
The period-end join leaks information about the real returns (the report moves them); on a noise target it carries
nothing, so the canary sees no gain ($-0.44\%$). The truncation test asks when each value became known and sees the
surprise one month early. Target encoding is built from the target: on a noise target it still leaks (the canary's
4.03\%), and recomputed with data truncated at a date it differs (the truncation test).
\end{solution}

\begin{solution}{exo:ml:validation:5}
The features locate the period, so shuffled folds would leak (neighbours are recognisable); and the test period's
feature distribution differs from training (drift, a new regime, a changed data source). It does not show that the model
is wrong or that the features predict returns: a feature can shift and stay useful.
\end{solution}

\begin{solution}{exo:ml:validation:6}
Whole-sample statistics use the test data: the scale and centre of each feature in the training rows depend on the
future. For stationary, well-behaved features the leak is small; for trending or drifting ones (volumes, prices,
counts) it is not, and it is a look-ahead however small. Fit the scaler inside each training fold, or use
cross-sectional ranks computed date by date.
\end{solution}

\begin{solution}{exo:ml:validation:7}
\texttt{announcement\_effect(0.04)}: the period-end join scores 5.47\% (1.27\% at 2\%), the filing-date join $-0.14\%$.
The explained variance from a known surprise is proportional to the square of the effect, so doubling it multiplies
the gain by about four; the model's own error adds the same small negative term in both cases.
\end{solution}

\begin{solution}{exo:ml:validation:8}
In each outer fold the procedure (search, choice, refit) uses only the outer training rows, so its fitted model is
independent of the outer test rows; its score there is an unbiased estimate of the generalisation score of what the
procedure produced from that much data, and so is the average over folds. The flat best takes a maximum over candidates
of scores measured on the same test rows, and by \cref{prop:ml:validation:max} its expectation exceeds every candidate's
true score.
\end{solution}

\begin{solution}{pb:ml:validation:1}
\begin{enumerate}
\item $-0.19\%$: nothing is predictable at decision time, so the best possible score is zero and a fitted model can only
lose to it.
\item The last training row's return is that of month 160: a gap keeps the training labels from ending inside the test.
\item Filing date: the latest surprise filed by the end of month $t$ (quarter ended at $t-1$ or before). Period end: the
surprise of the quarter ending at $t$, filed during $t+1$.
\item 2\% of return per unit of standardised surprise.
\item 8.97\%: the same-month industry return explains the industry factor's share of the variance, which is large.
\item 0.13\% on the whole sample, $-0.16\%$ on the training months.
\item 38\,080 against none.
\item 3.47\% and $-3.37\%$; persistent characteristics let deep trees recognise a stock in a given period and read its
neighbours' overlapping labels.
\item From $-0.60\%$ to $-0.27\%$; the best is 0.13 points above the median $-0.40\%$.
\item $-0.20\%$, against a median of $-0.19\%$; the two periods' scores correlate at $-0.16$.
\item Flat best $-0.62\%$, nested $-0.85\%$.
\item The maximum of noisy estimates overstates in expectation; the nested score is measured on data the choice never
saw.
\item Period-end join: truncation test (4.77 against 0). Target encoding: canary (4.03\%) and truncation test.
Screening: canary (0.17\%) and truncation test (5.13). Overlap: fold-overlap check (38\,080 against 0). Selection:
untouched holdout or nested search.
\item Each detector looks at one place (features, pipeline, splits, choice); a leak lives in one place.
\item The characteristics identify the period, so any shuffled split leaks and the test months differ from the training
months.
\item \emph{Named result}: on a target that cannot be predicted, the period-end join fakes 1.27\%, target encoding 8.97\%,
whole-sample screening 0.13\% ($-0.16\%$ honest), shuffled folds on overlapping labels 3.47\% ($-3.37\%$ purged), and
selecting the best seed on the test 0.13 points; truncation, canary, fold-overlap and holdout catch them in that order.
\item Features recomputed as of sample dates match; canary score within two standard errors of chance; no training
label overlaps a test label; adversarial AUC recorded; holdout untouched; trial count recorded.
\item None survives: each vanishes as soon as the data arrive in real time. Selection shows up as a smaller live score;
the others make live trading lose at once, which is why they are found late if paper trading is short.
\item As a finding with its measured effect, the fix, and every result it invalidates.
\item To estimate honestly how a procedure will do on data it has not seen, which requires that it really has not seen
them.
\end{enumerate}
\end{solution}

\begin{solution}{iq:ml:validation:1}
Information in training or features that would not be available at the decision time. Examples: fundamentals joined by
period end instead of filing date; revised macro data used as first released; universe of today's index members
(survivorship); features scaled with future statistics; overlapping labels across shuffled folds.
\iqlookfor{a precise definition tied to decision time, and concrete financial cases.}
\end{solution}

\begin{solution}{iq:ml:validation:2}
The best of twenty fold-averages is biased upwards (a maximum of noisy estimates). Report a \omterm{def:ml:validation:nested}{nested cross-validation} or a
walk-forward score of the whole procedure, and the number of settings tried.
\iqlookfor{selection bias and the nested or holdout remedy.}
\end{solution}

\begin{solution}{iq:ml:validation:3}
Recompute each feature on data truncated at a set of dates and compare with the stored feature (any difference is
look-ahead); run the pipeline on noise targets with the same structure (canary); check the knowledge times of inputs
against decision times.
\iqlookfor{black-box tests that do not rely on reading code.}
\end{solution}

\begin{solution}{iq:ml:validation:4}
A classifier trained to separate training from test rows; its AUC measures how different they are. High AUC warns that
the split's sides differ (drift) and that the features locate time, so shuffled splits leak; the features it relies on
are the ones to inspect.
\iqlookfor{what it measures and the two uses.}
\end{solution}

\begin{solution}{iq:ml:validation:5}
Encodings computed with same-period or future targets leak the target itself. Compute them from past periods only
(expanding or lagged means), inside each training fold, with smoothing toward the global mean.
\iqlookfor{lagging the statistic, and computing it inside the fold.}
\end{solution}

\begin{solution}{iq:ml:validation:6}
Run the \omterm{def:ml:validation:audit}{leakage audit} (truncation, canary, fold overlap): a leak shows there. Then check the search: trial count,
nested score, deflated statistics; a nested score near zero means \omterm{def:ml:why-financial-machine-learning-is-different:generalisation}{overfitting}. If both pass, compare feature
distributions and model errors live against backtest (\omterm{def:ml:validation:audit}{adversarial validation}, drift tests, chapter~12) and compute how
long the paper-trading period must be for the gap to be significant.
\iqlookfor{an ordered diagnosis with a test for each cause.}
\end{solution}
